% Einheit 2 – Terme zusammenfassen (bank/terme/e2.jsonl, zusammenbau v1.2)
\einheitenkopf*[e2][Terme zusammenfassen]{Terme zusammenfassen}
\setcounter{aufgabe}{5}


% test: terme-e2-k4-s11-v2, terme-e2-k5-s7-v2
\lbtest{Kannst du das schon? Dann weiter zu „Klammern auflösen“}{
\lbtt{a)}{Vereinfache den Term $10b - 5b^2 - 7b$ und berechne seinen Wert für $b = 2$. \quad \lbfeld}
\lbtt{b)}{$(-4x) \cdot 3y \cdot (-2) =$ \lbfeld}
}{a) $-14$ \quad b) $24xy$}


% erkennung: terme-e2-k1-s0-v1, terme-e2-k1-s0-v2, terme-e2-k1-s0-v3, terme-e2-k1-s0-v4
\begin{kbaufgabe}{Gleichartige Glieder erkennen}
\begin{kbblock}
\kbanweisung{Unterstreiche, was zusammengehört.}
\item[]\hspace*{-\leftmargin}\lbzelle{2}{\kbteil{a)}{$4y$, $7$, $9y$}{}}\hfill\lbzelle{2}{\kbteil{b)}{$6$, $2a$, $8$, $a$}{}}\par
\end{kbblock}
\begin{kbblock}
\item[]\hspace*{-\leftmargin}\lbzelle{2}{\kbteil{c)}{$5x^2$, $4x$, $3x^2$}{}}\hfill\lbzelle{2}{\kbteil{d)}{$7b$, $2c$, $5b$, $4$}{}}\par
\end{kbblock}
\end{kbaufgabe}


% erkennung: terme-e2-k2-s0-v1, terme-e2-k2-s0-v2, terme-e2-k2-s0-v3, terme-e2-k2-s0-v4
\begin{kbaufgabe}{Die Vorzahl eines Gliedes ablesen}
\begin{kbblock}
\kbanweisung{Welche Vorzahl hat …?}
\item[]\hspace*{-\leftmargin}\lbzelle{3}{\kbteil{a)}{$a$: \lbfeld}{}}\hfill\lbzelle{3}{\kbteil{b)}{$-b$: \lbfeld}{}}\hfill\lbzelle{3}{\kbteil{c)}{$1{,}5z$: \lbfeld}{}}\par
\end{kbblock}
\begin{kbblock}
\item[]\hspace*{-\leftmargin}\lbzelle{3}{\kbteil{d)}{$-4w$: \lbfeld}{}}\hfill\lbzelleleer{3}\hfill\lbzelleleer{3}\par
\end{kbblock}
\end{kbaufgabe}


% erkennung: terme-e2-k3-s0-v1, terme-e2-k3-s0-v2, terme-e2-k3-s0-v3, terme-e2-k3-s0-v4
\begin{kbaufgabe}{Die Glieder eines Terms mit ihrem Vorzeichen aufschreiben}
\begin{kbblock}
\kbanweisung{Schreibe die Glieder des Terms … mit ihrem Vorzeichen auf.}
\item[]\hspace*{-\leftmargin}\lbzelle{3}{\kbteil{a)}{$5a - 2b + 7$: \lbfeld}{}}\hfill\lbzelle{3}{\kbteil{b)}{$-3x + 6 - y$: \lbfeld}{}}\hfill\lbzelle{3}{\kbteil{c)}{$2m - 9 - 5k$: \lbfeld}{}}\par
\end{kbblock}
\begin{kbblock}
\item[]\hspace*{-\leftmargin}\lbzelle{3}{\kbteil{d)}{$-6 + 3c - 4d$: \lbfeld}{}}\hfill\lbzelleleer{3}\hfill\lbzelleleer{3}\par
\end{kbblock}
\end{kbaufgabe}


% leiter: terme-e2-k4-s1-v1, terme-e2-k4-s1-v2, terme-e2-k4-s2-v1, terme-e2-k4-s3-v1, terme-e2-k4-s4-v1, terme-e2-k4-s5-v1, terme-e2-k4-s6-v1, terme-e2-k4-s7-v1, terme-e2-k4-s8-v1, terme-e2-k4-s9-v1, terme-e2-k4-s10-v1
\begin{kbaufgabe}{Terme zusammenfassen}
\begin{kbblock}
\item[]\hspace*{-\leftmargin}\lbzelle{3}{\kbteil{a)}{$7x + 2x =$ \lbfeld}{}}\hfill\lbzelle{3}{\kbteil{b)}{$7x + 4x =$ \lbfeld}{}}\hfill\lbzelleleer{3}\par
\end{kbblock}
\begin{kbblock}
\item[]\hspace*{-\leftmargin}\begin{lbflucht}\makebox[1.6em][r]{c)\hspace{0.2em}} & $9x - 4x$ & $=$ & \lbfeld \\[\lbfluchtabstand] \makebox[1.6em][r]{d)\hspace{0.2em}} & $2x + 4x + 3x$ & $=$ & \lbfeld \\[\lbfluchtabstand] \makebox[1.6em][r]{e)\hspace{0.2em}} & $5x + 3 + 2x$ & $=$ & \lbfeld \\[\lbfluchtabstand] \makebox[1.6em][r]{f)\hspace{0.2em}} & $4a + 3b + 2a$ & $=$ & \lbfeld \\[\lbfluchtabstand] \makebox[1.6em][r]{g)\hspace{0.2em}} & $y + 4y$ & $=$ & \lbfeld \\[\lbfluchtabstand] \makebox[1.6em][r]{h)\hspace{0.2em}} & $2x - 7x$ & $=$ & \lbfeld \\[\lbfluchtabstand] \makebox[1.6em][r]{i)\hspace{0.2em}} & $-3x + 8x$ & $=$ & \lbfeld \\[\lbfluchtabstand] \makebox[1.6em][r]{j)\hspace{0.2em}} & $4x^2 + 5x + 2x^2$ & $=$ & \lbfeld \\[\lbfluchtabstand] \makebox[1.6em][r]{k)\hspace{0.2em}} & $1{,}5x + 3{,}5x$ & $=$ & \lbfeld\end{lbflucht}\par
\end{kbblock}
\end{kbaufgabe}


% pruefung: terme-e2-k4-s11-v1
\begin{kbaufgabe}{Wie in der Prüfung}
\begin{kbblock}
\kbteil{}{Vereinfache den Term $8a - 2a^2 - 5a$ und berechne seinen Wert für $a = 3$.}{P10 ’25 G}
\item[]\kbraum{2}
\end{kbblock}
\end{kbaufgabe}


% leiter: terme-e2-k5-s1-v1, terme-e2-k5-s1-v2, terme-e2-k5-s2-v1, terme-e2-k5-s3-v1, terme-e2-k5-s4-v1, terme-e2-k5-s5-v1, terme-e2-k5-s6-v1, terme-e2-k5-s7-v1
\begin{kbaufgabe}{Terme malnehmen}
\begin{kbblock}
\item[]\hspace*{-\leftmargin}\lbzelle{3}{\kbteil{a)}{$5 \cdot 2x =$ \lbfeld}{}}\hfill\lbzelle{3}{\kbteil{b)}{$5 \cdot 3x =$ \lbfeld}{}}\hfill\lbzelleleer{3}\par
\end{kbblock}
\begin{kbblock}
\item[]\hspace*{-\leftmargin}\begin{lbflucht}\makebox[1.6em][r]{c)\hspace{0.2em}} & $4x \cdot 5$ & $=$ & \lbfeld \\[\lbfluchtabstand] \makebox[1.6em][r]{d)\hspace{0.2em}} & $4y \cdot 2y$ & $=$ & \lbfeld \\[\lbfluchtabstand] \makebox[1.6em][r]{e)\hspace{0.2em}} & $(-4) \cdot 5x$ & $=$ & \lbfeld \\[\lbfluchtabstand] \makebox[1.6em][r]{f)\hspace{0.2em}} & $4a \cdot 5b$ & $=$ & \lbfeld \\[\lbfluchtabstand] \makebox[1.6em][r]{g)\hspace{0.2em}} & $2 \cdot 5x \cdot 3$ & $=$ & \lbfeld \\[\lbfluchtabstand] \makebox[1.6em][r]{h)\hspace{0.2em}} & $(-6a) \cdot 2b$ & $=$ & \lbfeld\end{lbflucht}\par
\end{kbblock}
\end{kbaufgabe}


% pflicht: terme-e2-k6-s1-v3, terme-e2-k6-s2-v2, terme-e2-k6-s3-v1, terme-e2-k6-s4-v1
\begin{lbkasten}
\begin{kbaufgabe}{Verstanden?}
\begin{kbblock}
\item[]\lbhinweis{Gemischt – erkenne selbst, was zu tun ist.}
\kbteil{a)}{Jana hat vier Terme zusammengefasst. Welche Ergebnisse können nicht stimmen? Begründe, ohne genau zu rechnen.}{}
\kbfrage{$4x + 5x = 9x^2$ \\ $6y - 11y = 5y$ \\ $2a + 7a = 9a$ \\ $3b \cdot 2b = 6b$}
\item[]\kbraum{2}
\end{kbblock}
\begin{kbblock}
\kbteil{b)}{Entscheide bei jeder Aussage, ob sie wahr oder falsch ist. Begründe, ohne genau zu rechnen.}{}
\kbfrage{(1) Zwei Glieder mit derselben Variablen und derselben Hochzahl kann man immer zusammenfassen. \\ (2) Zwei Glieder mit derselben Vorzahl kann man immer zusammenfassen. \\ (3) Beim Zusammenfassen entsteht nie eine neue Hochzahl.}
\item[]\kbraum{3}
\end{kbblock}
\begin{kbblock}
\kbteil{c)}{Die Figur zeigt ein Rechteck. Die Maße sind in cm angegeben. Schreibe einen Term für den Umfang. Fasse den Term zusammen.}{}
\kbzweispaltig[0.48]{\viereck[seiten={3x,2,3x,2}]{(0,0)}{(6,0)}{(6,2)}{(0,2)}}{\lbantwortrechts[1.2cm]{\kbfeld}}
\end{kbblock}
\begin{kbblock}
\kbteil{d)}{Ein Kiosk verkauft Brötchen zu je $p$ Euro. Am Montag verkauft er 45 Stück, am Dienstag 38 und am Mittwoch 52. Stelle einen Term für die Einnahmen der drei Tage auf und fasse ihn zusammen. Der Kiosk will in den drei Tagen mindestens 50 € mit Brötchen einnehmen. Schafft er das bei $p = 0{,}40$?}{}
\item[]\kbraum{3}
\end{kbblock}
\end{kbaufgabe}
\end{lbkasten}
